% =====================================================================
\chapter{Vocabulario, tokenización y modelos de lenguaje}
\label{tema:1}
% =====================================================================

Este tema fija los objetos sobre los que trabaja todo el curso: el vocabulario, las secuencias de tokens y la distribución de probabilidad que llamamos modelo de lenguaje. La presentación sigue la introducción del capítulo~3 de~\cite{jm3} y la clase~4 de~\cite{cs224n}.

\section{Vocabulario y secuencias}

\begin{definicion}[Vocabulario]
Un \emph{vocabulario} es un conjunto finito no vacío $\V$ cuyos elementos se llaman \emph{tokens}. Se escribe $K=\abs{\V}$ y se fija una biyección $\V\to\{0,1,\dots,K-1\}$ que asigna a cada token su \emph{índice}.
\end{definicion}

\begin{definicion}[Secuencias]
Para $T\ge 0$, $\V^{T}$ denota el conjunto de secuencias $x_{1:T}=(x_1,\dots,x_T)$ con $x_t\in\V$; $\V^{0}$ contiene solo la secuencia vacía. El conjunto de todas las secuencias finitas es
\[
\V^{*}=\bigcup_{T\ge 0}\V^{T}.
\]
\end{definicion}

\begin{proposicion}\label{prop:contar-secuencias}
Si $K=\abs{\V}\ge 2$, el número de secuencias de longitud a lo sumo $L$ es
\[
\sum_{T=0}^{L} K^{T}=\frac{K^{L+1}-1}{K-1}.
\]
En particular, $\V^{*}$ es infinito numerable.
\end{proposicion}

\section{Tokenización}

El texto se representa como una cadena de caracteres de un alfabeto $\Sigma$ (por ejemplo, Unicode). Un modelo de lenguaje no opera sobre $\Sigma^{*}$, sino sobre secuencias de tokens.

\begin{definicion}[Tokenizador]
Un \emph{tokenizador} es una aplicación $\tau\colon\Sigma^{*}\to\V^{*}$ junto con una aplicación de \emph{decodificación} $\delta\colon\V^{*}\to\Sigma^{*}$ tal que $\delta(\tau(s))=s$ para toda cadena $s$. La condición implica que $\tau$ es inyectiva: dos textos distintos nunca producen la misma secuencia de tokens.
\end{definicion}

\addtocounter{definicion}{1}%

La elección de $\V$ determina la granularidad. El \cref{tab:granularidad} resume las tres opciones habituales; los modelos actuales usan subpalabras obtenidas con el algoritmo BPE, que se estudia en el módulo~12.

\begin{table}[h]
\centering\small
\caption{Granularidad de la tokenización.}\label{tab:granularidad}
\begin{tabular}{@{}lccl@{}}
\toprule
Unidad & $K$ típico & Longitud de la secuencia & Limitación principal \\
\midrule
Carácter & $\sim 10^{2}$ & muy larga & secuencias costosas de procesar \\
Palabra & $>5\cdot 10^{5}$ & corta & palabras fuera del vocabulario \\
Subpalabra (BPE) & $3\cdot 10^{4}$–$2{,}6\cdot 10^{5}$ & intermedia & requiere entrenar el tokenizador \\
\bottomrule
\end{tabular}
\end{table}

\section{Modelos de lenguaje}

\begin{definicion}[Modelo de lenguaje]\label{def:ml}
Un \emph{modelo de lenguaje} sobre $\V$ es una distribución de probabilidad $P$ sobre $\V^{*}$, es decir, una función $P\colon\V^{*}\to[0,1]$ tal que
\[
\sum_{x\in\V^{*}}P(x)=1.
\]
\end{definicion}

La suma es una serie, porque $\V^{*}$ es infinito (\cref{prop:contar-secuencias}). Que una construcción concreta cumpla la condición de normalización no es automático: el \cref{tema:2} muestra que un modelo de bigramas sin símbolo de fin la incumple.

\begin{observacion}
En la práctica no se trabaja con $P(x)$ directamente, sino con las probabilidades condicionadas del token siguiente, $P(x_t\mid x_{1:t-1})$. Un gran modelo de lenguaje es una función con parámetros $\theta$ que recibe el contexto $x_{1:t-1}$ y devuelve un vector de $K$ números reales (los \emph{logits}); la función softmax (\cref{tema:4}) los convierte en una distribución sobre $\V$. La regla de la cadena (\cref{tema:2}) conecta ambas descripciones.
\end{observacion}

